\newpage%\新一页
\section{Assumptions and Justifications}
In response to the title of this article, the following hypotheses are proposed:
%1.	假设题目所给的数据真实可靠；
%注意：假设对整篇文章具有指导性，有时决定问题的难易。一定要注意假设的某种角度上的合理性，不能乱编，完全偏离事实或与题目要求相抵触。注意罗列要工整。

\begin{enumerate}[
	label = (\arabic*),
	itemindent = 0pt,
	labelindent = \parindent,
	labelwidth = 2em,
	labelsep = 5pt,
	leftmargin = 8em]
	\item[\textbf{假设1：}] 题面提供的DWTS官方数据表在Q0预处理后可视为内部一致；由\texttt{season$\times$week$\times$contestant}构成的键是Q1--Q4的唯一主线口径。
	\item[\textbf{假设2：}] Q0提取的淘汰/退赛标记能够正确识别当周退出选手集合；退出后的缺失分数不做插补，以避免引入未来信息。
	\item[\textbf{假设3：}] 观众投票不可直接观测；Q1将观众偏好作为潜变量，并由退出结果与评委信号共同约束。因此Q1输出解释为周内\emph{相对}人气强度，而非绝对票数。
	\item[\textbf{假设4：}] 在Q2--Q4评估替代机制时，我们固定赛季-周背景，并将Q1推断得到的观众强度后验视为可行偏好空间；不额外建模规则变化诱发的策略性行为改变。
	\item[\textbf{假设5：}] \texttt{src/mcm2026/config/config.yaml}中的配置参数（如\texttt{alpha}, \texttt{tau}, \texttt{outlier\_mults}）体现制作方偏好与压力强度设定；鲁棒性结论均在这些设定条件下报告。
\end{enumerate}



\section{Notations}
% 这部分不要过页（删掉此句话）
\begin{center}
	\begin{tabular}{cccc}
		\toprule[1pt] 
		\makebox[0.15\textwidth][c]{Symbol} & \makebox[0.3\textwidth][c]{Description} & \makebox[0.15\textwidth][c]{Symbol} & \makebox[0.3\textwidth][c]{Description} \\  
		\hline
		$s$ & 赛季索引 & $t$ & 赛季内周索引 \\
		$i$ & 选手（celebrity）索引 & $J_{s,t,i}$ & 评委信号（例如\texttt{judge\_score\_pct}） \\
		$p_{s,t,i}$ & 观众份额（Q1后验汇总） & $\alpha$ & 机制中评委-观众权重 \\
		$\tau$ & Q1软约束温度 & ESS ratio & 有效样本量占比（Q1不确定性） \\
		\bottomrule[1pt]
	\end{tabular}
\end{center}
\vspace{-0.5em}
\noindent{\small Note: Undefined variables are defined where they first appear.}


